\section{Causal upper sieve policies}\label{app:upper}

This appendix gives the upper-policy density and coefficient realization used in Section~\ref{sec:current-upper}; its finite savings are evaluated in Appendix~\ref{app:numeric}.

We construct the divisor coefficients from finitely many decisions
on prime prefixes. Prefixes are ordered by decreasing primes:
if $p_i=N^{t_i}$, then $t_1>\cdots>t_r>0$.
Write $S_j=t_1+\cdots+t_j$ and $S_0=0$.
A parent is obtained by deleting the last prime of a prefix.
Every decision at that parent is retained when smaller primes
are appended.

\subsection{Admissible caps and the parity of first exits}

For a nonempty subset $I\subseteq\{1,\ldots,r\}$ write
\begin{equation}\label{sv:selected-gate}
 b_I=\sum_{i\in I}t_i,\qquad
 G_I(t)=\max_{\substack{j\le r\\j\notin I}}
 \left(b_I+\sum_{\substack{i<j\\i\notin I}}t_i+2t_j\right),
 \qquad \max\varnothing=0.
\end{equation}
The full-square support condition is
\begin{equation}\label{sv:full-gate}
 G_{\mathrm{full}}(t)=\max_{1\le j\le r}(S_{j-1}+2t_j).
\end{equation}
The candidate selected-product endpoint is
\begin{equation}\label{sv:tent}
 v(b)=\min\left\{\frac{1+b}{2},\frac{57}{64}-\frac{3b}{2}\right\}.
\end{equation}
A selected branch on a compact prefix region is admitted only when
its constant cap $\lambda$ satisfies, with a fixed positive margin,
\begin{equation}\label{sv:branch-licence}
 \max\{1/2,G_{\mathrm{full}}(t),G_I(t)\}
       <\lambda<v(b_I)
\end{equation}
throughout that region.  The separate three-slot branch with endpoint
$\Lambda_{\rm cap}-3t_3/8$ is subject to the whole-box factorization conditions
of Appendix~\ref{app:coefficients}.  A maximum of candidate endpoints
is always taken over branches satisfying their support conditions.

An ordinary upper residual tree at total exponent $\lambda$, after
a prefix of rank $r$, applies the cube tests
\begin{equation}\label{sv:upper-cubes}
 S_{j-1}+3t_j<\lambda\qquad(j>r,\ j\text{ odd}).
\end{equation}
If $r$ is even, these are the ordinary upper residual tests.  If $r$
is odd, they are the ordinary lower residual tests.  The earlier
membership decisions are retained.  At an even prefix a change of cap
is made after admitting the prefix; it introduces no exclusion at that
rank.  At an odd prefix a cap may be chosen before its own odd test,
provided that the decision depends only on that prefix and the
previously fixed policy.  Either operation leaves every first
exclusion at odd full rank.

\begin{lemma}[First-exit identity]\label{sv:first-exit}
Let $\mathcal T$ be a prefix-closed collection of squarefree products
from a finite prime set $\mathcal P$, containing the empty product.
Suppose that all first excluded prefixes have odd rank.  For every
multiplicative $\xi$ with $0\le \xi(p)\le1$,
\begin{equation}\label{sv:upper-exit}
 \sum_{d\in\mathcal T}\mu(d)\xi(d)
 =\prod_{p\in\mathcal P}(1-\xi(p))
 +\sum_{e\text{ first exit}}\xi(e)
           \prod_{\substack{p\in\mathcal P\\p<p_{\mathrm{last}}(e)}}
                            (1-\xi(p)).
\end{equation}
In particular $\sum_{d\mid n,\ d\in\mathcal T}\mu(d)$ majorizes
$\one_{(n,\prod_{p\in\mathcal P}p)=1}$ for every integer $n$.
\end{lemma}

\begin{proof}
Each excluded divisor has a unique first excluded prefix $e$; its
remaining prime factors are smaller than $p_{\mathrm{last}}(e)$.
Summing the complete continuation gives
$\mu(e)\xi(e)\prod_{p<p_{\mathrm{last}}(e)}(1-\xi(p))$.
Subtract these disjoint continuations from the complete Euler product.
Since $\mu(e)=-1$, the sign in \eqref{sv:upper-exit} is positive.
Take $\xi(p)=\one_{p\mid n}$ for the last assertion.  The same
finite identity applies at $n=0$ and at negative integers.
\end{proof}

The square tests persist through descendants.  A retained odd cube
implies its own square.  The following even square follows from
ordering:
\begin{equation}\label{sv:descendant-square}
 S_{j-1}+2t_j\le S_{j-2}+3t_{j-1}<\lambda.
\end{equation}
After all selected slots have appeared, the expression in
\eqref{sv:selected-gate} for any appended slot equals its full square.
Consequently \eqref{sv:branch-licence}, checked at the selected prefix,
continues to control the complete descendant vector.

\subsection{The cap changes and their density differences}

Set
\[
 \rho_0=\frac{19}{96},\qquad \mu_0=\frac{25}{128},\qquad
 \Theta=\frac{299}{500},\qquad \Lambda_{\rm cap}=\frac{153}{256}.
\]
The first prime is less than $N^{\rho_0}$.  The three-slot endpoint
below $t_1=\mu_0$ is the envelope of admissible branches of
\begin{equation}\label{sv:third-envelope}
 v(t_1),\ v(t_1+t_2),\ v(t_1+t_3),\ v(t_2+t_3),\ v(S_3),\
                         \Lambda_{\rm cap}-\frac{3t_3}{8};
\end{equation}
above $\mu_0$ the initial three-slot endpoint is $v(t_1)$.
Every branch in \eqref{sv:third-envelope} retains its individual
factorization conditions.  In particular the displayed list is not
an unconditional permission to select any nonleading pair.

At rank four the selected candidates are $M=v(b_I)$ with
$4\in I$, satisfying \eqref{sv:branch-licence} and exceeding the
inherited third cap.  Such a change is made only on prefixes that
already pass the third cube.  Membership through rank four is
unchanged and subsequent odd cubes use the new cap.
The limiting density difference on a region $\mathcal C$ is
\begin{equation}\label{sv:even-gain}
 4\int_{\mathcal C}
 \frac{\overline F((L-S_r)/t_r)
                   -\overline F((M-S_r)/t_r)}
      {t_1\cdots t_{r-1}t_r^2}\,dt_1\cdots dt_r,
 \qquad \overline F=\frac{F}{2e^\gamma},
\end{equation}
where $r=4$, $L$ is the inherited cap and $M$ the selected endpoint.
Indeed an even prefix $P_r$ contributes
$g_N(P_r)V_N^+(N^\lambda/P_r,p_r)$ to its complete subtree.
Subtracting the two subtrees gives \eqref{sv:even-gain}; it is an
identity of raw densities before taking an asymptotic limit.

At rank five a candidate $v(b_I)$, $5\in I$, is compared with the
inherited cap and the fifth cube $S_4+3t_5$.  Its cap is fixed before
that cube.  Thus a fifth prefix rejected at the smaller cap may be
retained at the larger one; all earlier decisions remain fixed.
The negative fifth subtree has the ordinary lower residual, and its
saving is
\begin{equation}\label{sv:odd-gain}
 4\int_{\mathcal C}
 \frac{\overline f((M-S_5)/t_5)
                   -\overline f((L-S_5)/t_5)}
      {t_1t_2t_3t_4t_5^2}\,dt_1\cdots dt_5,
 \qquad \overline f=\frac{f}{2e^\gamma},
\end{equation}
with the zero extension of $\overline f$ below $2$.
The inherited first-five cap is bounded above by
\begin{equation}\label{sv:fifth-envelope}
 R_5(t)=\max\left\{\Lambda_{\rm cap}-\frac{3t_3}{8},\
         v(b_I):\varnothing\ne I\subseteq\{1,\ldots,5\}\right\}.
\end{equation}
Using an upper bound for the inherited cap in a positive saving is
conservative.  It does not change the physical cap or its density.

The rank-six changes are even-prefix changes of the form
\eqref{sv:even-gain}.  They are applied after preserving the third
and fifth admission tests.  In particular
$L_3-S_6>2t_5-t_6>t_6$ on a region whose fifth cube is controlled by
$L_3$.  A convex piecewise-affine majorant $U_5\ge R_5$ can be obtained by
choosing one defining affine upper plane for each tent in
\eqref{sv:fifth-envelope}.  For a new concave endpoint $M$, put
$d=M-U_5$ and restrict the region to $M-S_6\le3t_6$.
Write $L$ for the inherited cap.  Where $d>0$, the retained fifth
cube gives
\[
 1<\frac{L-S_6}{t_6}\le\frac{U_5-S_6}{t_6}
       <\frac{M-S_6}{t_6}\le3.
\]
Thus $\overline F(s)=1/s$ throughout the interval of comparison, and
\[
 \overline F((L-S_6)/t_6)-\overline F((M-S_6)/t_6)
 \ge \frac{t_6(M-U_5)}{(U_5-S_6)(M-S_6)}.
\]
The cap is changed only where $d>0$, giving the saving lower bound
\begin{equation}\label{sv:sixth-lower}
 \frac{4d_+}{R_oR_n\prod_{i=1}^6t_i},
\end{equation}
where $R_n\ge M-S_6$ and
$R_o\ge\min\{U_5-S_6,M-S_6\}>0$ are whole-region upper bounds.
The positive part is taken before a lower bound for the positive
prime-density measure is extracted.  For a rectangular region,
concavity gives
$\int d_+/\operatorname{Vol}(\mathcal C)
\ge 2^{-6}\sum_{v\text{ vertex}}d(v)$.

Two additional changes occupy disjoint parts of $\mu_0<x<\rho_0$.
For the rank-two change, put $a(x)=2v(x)-1=25/32-3x$ and use
\[
 a(x)<y<x,\qquad L=v(x),\quad M=v(y).
\]
The support conditions follow from
$2x,x+2y,2x+y< v(x)<v(y)$; the selected slot is $2$.
The saving is the $r=2$ instance of \eqref{sv:even-gain}.
For the rank-three change, put $t=y+z$ and impose
\[
 a(x)<t<x,\quad z=tw/2,\quad 5/9\le w<1,
 \qquad L=v(x),\quad M=v(t).
\]
Here $y<3/16$, whereas the rank-two change has $y\ge3/16$.
Selecting slots $\{2,3\}$ gives the extra support condition $2x+t<M$.
The third node is retained and its lower residual is changed, giving
the $r=3$ analogue of \eqref{sv:odd-gain}.  Its arguments lie in
$[2,4]$, where $\overline f(s)=\log(s-1)/s$.

At rank eight a selected endpoint $M$ is compared with a valid
upper bound $R_6$ for the inherited sixth cap.  The retained seventh
cube and the additional condition $M-S_8<3t_8$ give
\begin{equation}\label{sv:eighth-range}
 1<\frac{\lambda_6-S_8}{t_8}
 \le\frac{R_6-S_8}{t_8}<\frac{M-S_8}{t_8}<3.
\end{equation}
Consequently \eqref{sv:even-gain} is bounded below by the integral of
\begin{equation}\label{sv:eighth-density}
 \frac{4(M-R_6)}{(R_6-S_8)(M-S_8)\prod_{i=1}^8t_i}.
\end{equation}
Each chosen region must retain every full and selected-complement
square, as well as all ancestor tests.  Regions acted upon at one
stage have disjoint interiors.  Changes at different stages are
combined by the exact telescope
$V_3-V_6=(V_3-V_4)+(V_4-V_5)+(V_5-V_6)$.
For each density difference we use one lower enclosure.

\subsection{Literal ordinary densities and finite realization}

Let $g_N(p)=\one_{p>2,\ p\nmid N}/(p-1)$ and
$V_N(y)=\prod_{p<y}(1-g_N(p))$.  The ordinary cubic Rosser
coefficients satisfy, uniformly on compact subsets of the stated
ranges,
\begin{equation}\label{sv:literal-density}
 \frac{V_N^+(N^\beta,N^t)}{V_N(N^t)}=F(\beta/t)+o(1)
       \quad(\beta/t>1),\qquad
 \frac{V_N^-(N^\beta,N^t)}{V_N(N^t)}=f(\beta/t)+o(1)
       \quad(\beta/t\ge2).
\end{equation}
Here the densities are those of the literal coefficients, not merely
upper and lower bounds for an arbitrary sifted sequence.

To see the uniformity in $N$, compare a normalized fixed-rank exit
with the corresponding exit for $g_0(p)=\one_{p>2}/(p-1)$.
A surviving prefix gains the multiplier
\[
 \prod_{\substack{p_{\rm last}\le p<y\\p\mid N}}(1-g_0(p))\le1;
\]
other prefixes disappear.  A rank-$j$ exit has last prime at least
$N^{\beta_{\min}/(3\cdot2^{j-1})}$ on a positive compact parameter
range.  Deleting prime divisors of $N$ therefore changes any
fixed-rank integral by $o(1)$ uniformly.  The positive Rosser exit
series converges uniformly on compact parameter ranges, and the
raw one-sided bounds give the opposite side of a finite-rank
sandwich; see \cite[Proposition~11.8 and Lemma~14.7]{Suzuki}.
First choose the rank and then let $N$ increase.  This proves
\eqref{sv:literal-density}.  For $1<s\le3$ the identity at the first
cube reduces the upper calculation to cutoff $N^{\beta/3}$, giving
$F(s)=2e^\gamma/s$.

The normalization follows from Mertens' formula.  In particular the
exact ratio to $2C_N\prod_{p<z}(1-1/p)$ is
\begin{equation}\label{sv:euler-ratio}
 \frac{V_N(z)}{2C_N\prod_{p<z}(1-1/p)}
 =\prod_{p\ge z}\left(1-\frac1{(p-1)^2}\right)^{-1}
  \prod_{\substack{p\mid N\\p\ge z}}\left(1-\frac1{p-1}\right).
\end{equation}
Its logarithm is $O((1+\log N/\log z)/(z-2))$.
Thus $V_N(N^t)\sim2e^{-\gamma}C_N/(t\log N)$ uniformly in $N$.
The factor $4$ in \eqref{sv:even-gain} and \eqref{sv:odd-gain}
is exactly $2e^{-\gamma}\cdot2e^\gamma$.

The finite cell data consist of rational regions, their ranks,
selected subsets, endpoints, and parent labels. Every active region
has positive margins in its admission and support inequalities;
an inherited region keeps its parent's decision. The density is
the sum of the exact subtree differences. For a positive saving,
any selected subfamily gives a lower bound. For a negative cost,
the upper bound must include the whole retained domain.

\begin{proposition}[Finite-policy realization]\label{sv:finite-realization}
Fix such a finite admissible policy and a positive error $\epsilon$.
Assume the coefficient factorization conditions of
Appendix~\ref{app:coefficients} on its selected branches.
There is a fixed finite refinement and a bounded squarefree
coefficient $\rho_N^{\rm up}$, supported on $(q,N)=1$, with the following
properties.  Its divisor sum is a nonnegative rough-indicator
majorant.  The signed distribution estimate \eqref{eq:current-PP}
applies to that same coefficient.  Its density differs, in the
required upper direction, from the ideal policy density by at most
$\epsilon C_N/\log N+o(C_N/\log N)$.
\end{proposition}

\begin{proof}
First truncate any small-last-prime tail.  At the third activation
one may use, with $u=\rho_0/2$ and $0<h<u$,
\begin{equation}\label{sv:activation-tail}
 T_3(h)\le\frac{18h}{(u-h)^2}
   \left(\log^2\frac{\rho_0}{h}
                   +2\log\frac{\rho_0}{h}+2\right)=o(1).
\end{equation}
This follows from $\overline F(s)-\overline f(s)\le9/s^2$
for $s\ge3$, the residual lower bound $\lambda-t_1-t_2\ge u$,
and the ordered reciprocal-prime triangle of mass
$\tfrac12\log^2(\rho_0/t_3)$.  Higher selected changes have
positive last-prime cutoffs.  Proper affine walls have arbitrarily
small strip volume.  Identities and open inherit regions are
retained.  On the remaining compact pieces choose coarse caps with
strict margins and refine them in increasing prefix order.  No
ancestor decision may depend on an appended prime.

The third endpoint is piecewise affine and Lipschitz.  For boxes
of width $\chi<e/10$, caps of the form
$\lambda_{\mathscr B}=\Gamma_{\mathscr B}-2e$ approximate it with error at most
$4e+6\chi$.  Near a three-slot branch with endpoint
$W=\Lambda_{\rm cap}-3t_3/8$, an auxiliary factorization cover may exceed
the physical cap.  If $\inf_{\mathscr B} A\ge\lambda_{\mathscr B}+e/2$, where
$A=(1+t_1)/2$, take $\ell_{\rm aux}=\lambda_{\mathscr B}+e/4$.
Otherwise $\sup_{\mathscr B} A<\lambda_{\mathscr B}+11e/20$, and
$\ell_{\rm aux}=\lambda_{\mathscr B}+e$ lies at least $9e/20$ above the
necessary support condition and $e$ below $W$.  The other initialization square
has margin $3/256$.  The auxiliary cover gives a factorization for the coefficient;
it never replaces the physical cap in a density or residual level.

Choose a presieve $z_0=N^{\eta_0}$.  At even rough rank use outer
endpoint tests and at odd rough rank use inner tests, retaining all
ancestor labels.  This gives $a_N(d)\ge\lambda_N(d)$ for the
raw upper coefficients.  If $B^-\le I_0\le B^+$ are the
dimension-one small-prime beta polynomials, attach $B^+$ to
positive $a_N(d)$ and $B^-$ to negative $a_N(d)$.  Term by term,
$a_N(d)B_{\rm sign}\ge a_N(d)I_0$.  Unique small/rough
factorization defines the resulting $\rho_N^{\rm up}$ and preserves
squarefreeness and boundedness.

The density direction needs an additional identity.  Projecting
\eqref{sv:upper-exit} to primes at least $z_0$ gives
\begin{equation}\label{sv:upper-projection}
 V_N(z_0)\sum_{d\text{ rough}}\lambda_N(d)g_N(d)
       \le\sum_{d\text{ full}}\lambda_N(d)g_N(d).
\end{equation}
Only positive exits with last prime at least $z_0$ survive on the
left.  For a fixed coarse policy the absolute beta and fine-box
errors are
$O(e^{-c/\eta}\eta_0^{-2}+C_{\rm policy}\tau\eta_0^{-5})$
relative to the root Euler product.  Choose $\eta,\eta_0$ and
then $\tau$ so that these errors and the truncated pieces fit
within $\epsilon$.

All these choices precede $N$.  The fine rank is bounded by
$J_0>\theta_*/\eta_0+2$; support inflation can be kept below
$1/1000$, giving $\theta_*=599/1000$.
Finally choose the fixed moment order and complete-scale trim so
that $\theta_*/(1-\kappa_0)<3/5$ and all local inequalities have
strict margins.  Appendix~\ref{app:coefficients} then gives
\eqref{eq:current-PP} for this very coefficient.  The pointwise
bound follows from Lemma~\ref{sv:first-exit}; the density bound
follows from \eqref{sv:upper-projection} and the absolute errors.
\end{proof}

\section{Lower policies and finite feedback}\label{app:lower}

This appendix supplies the lower-policy realizations and positive marked transition inequalities used in Section~\ref{sec:current-lower}. Appendix~\ref{app:numeric} evaluates the corresponding finite maps and integrals.

\subsection{The ordinary lower tree and fourth-prefix restoration}

Put ${a_{\rm s}}=417/5000$, ${h_{\rm s}}=19/125$ and $\nu=249/448$.
For either $w={a_{\rm s}}$ or $w={h_{\rm s}}$, begin with the ordinary lower support
at total exponent $\nu_\delta=\nu-\delta$, where
$0<\delta$ is fixed and $2{h_{\rm s}}<\nu_\delta$.
Its first square $2t_1<\nu_\delta$ is automatic below $w$;
its subsequent tests are
\begin{equation}\label{sv:lower-cubes}
 S_{j-1}+3t_j<\nu_\delta\qquad(j\text{ even}).
\end{equation}
At an admitted third prefix, freeze $\lambda_3\ge\nu_\delta$
for future even cubes.  An ordinary fourth change is permitted only
after $S_3+3t_4<\lambda_3$; it freezes
$\lambda_4\ge\lambda_3$.  The fifth prefix is then admitted, and
may freeze $\lambda_5\ge\lambda_4$ for future even tests.
Every active selected cap is below its admissible endpoint and above
its parent and all support conditions \eqref{sv:selected-gate}--\eqref{sv:full-gate}.
For the lower starting families the selected endpoint is
\begin{equation}\label{sv:lower-endpoint}
 v_-(b)=\min\left\{\frac{1+b}{2},
 \frac{32-34b}{39},\frac{57}{64}-\frac{3b}{2}\right\}.
\end{equation}
The selected branch is kept through descendants; the residual size
ratio of a marked state, introduced below, is separate from the
local residue exponent in the distribution estimate.

There is also a fourth-prefix restoration.  On an otherwise first
excluded fourth prefix choose $\lambda_4$ before its own cube,
subject to
\begin{equation}\label{sv:restore-gates}
 \lambda_3\le S_3+3t_4<\lambda_4,
 \qquad \max\{G_{\rm full},G_I\}<\lambda_4<v_-(b_I).
\end{equation}
Attach the ordinary lower residual below $p_4$ at level
$N^{\lambda_4}/P_4$.  Since
$\lambda_4-S_4>2t_4$, its first square cannot fail.
Its later first exits have even residual rank, hence even full
rank.  The restored tree remains a lower tree.  Its density gain is
exactly
\begin{equation}\label{sv:restoration-density}
 4\int_{\mathcal B}
 \frac{\overline f((\lambda_4-S_4)/t_4)}
      {t_1t_2t_3t_4^2}\,dt_1dt_2dt_3dt_4.
\end{equation}
There was no retained residual at such a first exit, so no parent
residual term is subtracted in \eqref{sv:restoration-density}.

Define a finite rational family of restoration regions as follows. Let
$r_0=(\nu-{h_{\rm s}})/9$, $n\in\{400,800\}$, and
$\mathscr B=\prod_{i=1}^4[\ell_i,\ell_i+1/n]$ with
$r_0\le\ell_4<\ell_3<\ell_2<\ell_1<{h_{\rm s}}$ and
$u_1+3u_2<\nu$, where $u_i=\ell_i+1/n$.
Let $U_3$ be the maximum of $\nu$, the upper bounds of all
three-slot selected endpoints, and $\Lambda_{\rm cap}-3\ell_3/8$.
It bounds the inherited third cap from above.
With $D=\operatorname{lcm}(448,256,78n,64n,8n)$, retain a mask
$I$ containing $4$ only if
\[
 c_I=\inf_{\mathscr B} v_-(b_I)-D^{-1}
       >\max\{U_3,\sup_{\mathscr B} G_{\rm full},\sup_{\mathscr B} G_I\}.
\]
Choose the largest retained $c_I$; write it as $c$.
For $A=t_1+t_2+t_3+3t_4$, clip the box to
\[
 \max(A_-,U_3+D^{-1})<A<\min(A_+,c-D^{-1}),
\]
and split this interval into finitely many equal rational layers.
Every layer satisfies \eqref{sv:restore-gates} with positive
margins and lies outside previously retained fourth prefixes.
On the clipped region,
\[
 2<\frac{c-S_4}{t_4}
       <2+\frac{9(\Lambda_{\rm cap}-\nu)}{\nu-{h_{\rm s}}}
       =2+\frac{84375}{90452}<3.
\]
Thus $\overline f(s)=\log(s-1)/s$ is increasing throughout.
If $A<a_j$ on a layer and $u_4$ is the fourth upper endpoint, use
$s_j=2+(c-a_j)/u_4$ in the lower bound
$4\overline f(s_j)\operatorname{Vol}(E_j)/(u_1u_2u_3u_4^2)$.
The exact volume of $\mathscr B\cap\{A<t\}$ is
\begin{equation}\label{sv:slab-volume}
 \frac1{72}\sum_{e\in\{0,1\}^4}(-1)^{|e|}
 \left(t-A_--\frac{e_1+e_2+e_3+3e_4}{n}\right)_+^4.
\end{equation}
This is the inclusion--exclusion formula for weighted box widths
$(1,1,1,3)/n$.  Differences of \eqref{sv:slab-volume} give
the layer volumes exactly.

\subsection{The lower coefficient and its original prime sequence}

For a prefix-closed tree with even first exits the counterpart of
\eqref{sv:upper-exit} is
\begin{equation}\label{sv:lower-exit}
 \sum_{d\in\mathcal T}\mu(d)\xi(d)
 =\prod_p(1-\xi(p))
 -\sum_{e\text{ first exit}}\xi(e)
               \prod_{p<p_{\rm last}(e)}(1-\xi(p)).
\end{equation}
Its proof is identical, using $\mu(e)=1$.
Taking $\xi(p)=\one_{p\mid n}$ gives a rough-indicator
minorant.  With $\xi=g_N$ and presieve $z_0$ it gives
\begin{equation}\label{sv:lower-projection}
 V_N(z_0)R_{\rm rough}-V_{\rm full}
 =\sum_{\substack{e\text{ full even exit}\\p_{\rm last}(e)<z_0}}
                      g_N(e)V_N(p_{\rm last}(e))\ge0.
\end{equation}
This is the density direction required for lower coefficient realization.

Round positive even coefficients inward and negative odd coefficients
outward, preserving the prefix decisions.  Attach beta-minus to a positive
rounded coefficient and beta-plus to a negative one.  Term by term
$a_{\rm rnd}(d)B_{\rm sign}\le a_{\rm rnd}(d)I_0$.
Consequently the physical coefficient satisfies
\begin{equation}\label{sv:physical-lower}
 \sum_{q\mid n}\rho_N^{\rm low}(q)\le
       \one_{(n,P_N(N^w))=1},\qquad
 \sum_q\frac{\rho_N^{\rm low}(q)}{\varphi(q)}
       \ge V_{\rm full}-\text{absolute density error}.
\end{equation}
The second inequality uses \eqref{sv:lower-projection} and the
absolute beta and box estimates.  It does not follow merely from
coefficientwise order, and beta-minus need not be nonnegative.
The sharp-prime consequence in Appendix~\ref{app:cross-review}, applied
to this same $\rho_N^{\rm low}$, turns \eqref{sv:physical-lower} into the
ordinary lower estimate for $\{N-p:p<N,\ p\text{ prime}\}$.

For an admissible lower prefix, the third-, fourth-, and fifth-stage
density changes are, respectively, upper parent minus upper child,
lower child minus lower parent, and upper parent minus upper child.
Evaluate each by~\eqref{sv:literal-density}.
At ranks $r=3,4,5$, the difference of the ordinary upper and lower
densities cancels the apparent inverse-square singularity in the last
coordinate, leaving a tail
$O(\epsilon(1+\log^{r-1}({h_{\rm s}}/\epsilon)))$.
The remaining compact regions admit one common finite refinement
that preserves every inherited decision.

\subsection{Finite feedback on positive marked families}\label{sv:finite-feedback}

The normalized states are estimates averaged over a nonnegative finite
marked family $\sigma(B)$ with a common total level $Q=N^\lambda$.
Put
\[
 L=\lambda-\log_NB>0,\qquad
 \Theta_Q=4\li(N)\sum_B
       \frac{\sigma(B)C_{BN}}{\varphi(B)\log(Q/B)},\qquad
 \alpha=\frac{1-\log_NB}{L}\ge2.
\]
For a singleton child of exponent $Lu$, or a pair of exponents $Lu,Lv$,
the residual level and parameter are, respectively,
\begin{equation}\label{sv:child-state}
 L'=L(1-u),\quad
 \alpha'=1+\frac{\alpha-1}{1-u};\qquad
 L'=L(1-u-v),\quad
 \alpha'=1+\frac{\alpha-1}{1-u-v}.
\end{equation}
The total $Q$, selected slots and prime ceiling are retained.
In particular $\alpha$ increases; the child is not reset to $\alpha=2$.

An admitted marked family here means a finite nonnegative combination
of labelled prime-box indicators, with fixed positive lower bounds on
all marked prime exponents and on $L$. In the classical case its
complete modulus level is below $N^{1/2}$ by a fixed power. In the
enhanced case $1/2<\lambda<153/256$ with fixed margins; it retains all full and selected-complement square
conditions and one factorization branch from
Appendix~\ref{app:coefficients}. It also retains a number $\chi>0$
such that adjoining primes of exponent below $\chi$ preserves that
branch and its support conditions. The number $\chi$ is no larger than the least
specified exponent. If a branch uses a bound $c$ for a third slot
not yet specified, it also satisfies $\chi\le c$.
Every enhanced ordinary node at argument $s$ has physical cutoff
$N^{L/s}$ below $N^\chi$. All these restrictions are imposed with
fixed margins on each box before endpoint limits are taken. The
estimates concern the marked average, rather than an individual fixed
prime prefix.

The child-to-parent measures are
$du/[u(1-u)]$ and $du\,dv/[uv(1-u-v)]$.
They are positive.  To specify the one-step switched cost, define
\begin{equation}\label{sv:one-step-mass}
 I(s,s_0,\alpha)=
 \int_{1/s_0<u_1<u_2<u_3<1/s}
 \omega\!\left(\frac{\alpha-u_1-u_2-u_3}{u_2}\right)
                    \frac{du_1\,du_2\,du_3}{u_1u_2^2u_3}.
\end{equation}
The one-step inequality gives the candidate
\begin{align}\label{sv:one-step-operator}
 \mathcal W_1[H,h](s;s_0,\alpha)
 ={}&A_+(s)-A_+(s_0)+H(s_0)-\frac K8 I(s,s_0,\alpha)\notag\\
 &+\frac12\int_{1/s_0}^{1/s}
 \frac{A_-(s_0(1-u))+h(s_0(1-u))}{u(1-u)}\,du,
 \qquad 2\le s<s_0\le5.
\end{align}
Here $A_+(s)=sF(s)/(2e^\gamma)$ and
$A_-(s)=sf(s)/(2e^\gamma)$, with the zero extension below
$2$.  The physical singleton child cutoff is $N^{L/s_0}$ and its
argument is $s_0(1-u)$, not the argument at a different total level.
The positive switched family is
$Bq_1q_2q_3n\le N$ with
$N^{L/s_0}\le q_1<q_2<q_3<N^{L/s}$ and $P^-(n)\ge q_2$.
Its unsifted mass, in the parent normalization, is $I/4$; the
factor $1/2$ in the pointwise weighted inequality gives the cost
$KI/8$ in \eqref{sv:one-step-operator}.  All its prime factors
have fixed positive exponents.  Its unit quotient has a fixed power
saving because $\alpha-3/s\ge1/2$.

A uniform bound usable in \eqref{sv:one-step-mass} is obtained by
putting $c_\omega=70893/125000$ and
\begin{equation}\label{sv:omega-envelope}
 \omega_{\rm env}(u)=
 \begin{cases}
 1,&u<1,\\
 \max\{1/u,c_\omega\},&u\ge1.
 \end{cases}
\end{equation}
For $v\ge u\ge1$ one has $\omega(v)\le\omega_{\rm env}(u)$.
On $[1,2]$ this follows from $\omega(v)=1/v$; on $[2,3]$
maximize $1+\log(v-1)-c_\omega v$.  The maximum is
$-\log c_\omega-c_\omega<0$, since the exact rational inequality
$c_\omega\sum_{j=0}^8c_\omega^j/j!>1$ implies
$c_\omega e^{c_\omega}>1$.  The delay equation propagates the
bound for $v\ge3$.  Thus replacing the integrand by
$\omega_{\rm env}((2-u_1-u_2-u_3)/u_2)$ bounds the cost for every
$\alpha\ge2$.  The branch below $1$ is a conservative bound for
the switched estimates that use it; unit quotient terms are counted
separately when their domains have nonnegligible mass.

For a two-step action write $a_0=1/s_0$, $a_i=1/\phi_i$ and $b=1/s$.
The parameter conditions are
\begin{gather}\label{sv:two-step-gates}
 2<s<\phi_3<\phi_2<\phi_1<s_0\le5,\qquad
 a_3+2a_1<1,\quad3a_1<1,\quad a_2+2a_1<1,\notag\\
 a_0<(1-a_2)/2<1/3<a_2<1/2,\qquad
 s_0(1-b)>2,\quad s_0-1\le4,\quad s_0(2-2a_2)>2.
\end{gather}
Its three singleton regions are
\begin{equation}\label{sv:singleton-domains}
 (a_0,b),\qquad(a_0,a_2),\qquad(a_0,a_3),
 \quad \rho(u)=s_0(1-u),
 \quad\text{physical cutoff }N^{La_0}.
\end{equation}
The regions on which the ordinary pair estimate and its feedback are used are
\[
\begin{array}{c|c|c}
 &\text{region}&\text{child argument}\\ \hline
 D_5^+&a_0<u<v<a_2,\ v+2u<1&s_0(1-u-v)\\
 D_6&(a_0,a_1)\times(a_2,a_3)&s_0(1-u-v)\\
 D_7&a_0<u<v<a_1&(1-u-v)/u\\
 D_8&(a_0,a_1)\times(a_1,a_2)&(1-u-v)/u,
\end{array}
\]
For the pointwise weighted inequality the fifth term has the complete
mother region
\begin{equation}\label{sv:complete-fifth-region}
 D_5=\{a_0<u<v<a_2\}
     =D_5^+\sqcup D_5^-,\qquad
 D_5^-=D_5\cap\{v+2u\ge1\}.
\end{equation}
Thus $\Gamma_5=\Gamma_5^++\Gamma_5^-$; it is not restricted to
$D_5^+$ in the combinatorial inequality. The ordinary estimate and
child defect are used only on $D_5^+$. On $D_5^-$ the complete count
is retained and bounded by the positive product estimate.
The test $2v<1$ holds throughout $D_5$, and the other three pair
regions satisfy both $2v<1$ and $v+2u<1$ by
\eqref{sv:two-step-gates}.
If $J_2^0$ denotes the ordinary and positive switched base for the
chosen admissible parameter tuple, the feedback added to that base is
\begin{equation}\label{sv:two-step-feedback}
 \frac15\left\{4H(s_0)+H(\phi_1)
 +\sum_{c\in\{b,a_2,a_3\}}\int_{a_0}^{c}
       \frac{h(s_0(1-u))}{u(1-u)}\,du
 +\sum_{i=5^+,6,7,8}\iint_{D_i}
       \frac{H(\rho_i(u,v))}{uv(1-u-v)}\,du\,dv\right\}.
\end{equation}
We now define this base and prove the sign of every contribution.

\paragraph{The twenty-one-term inequality.}
For a marked state let $R=Q/B$ and, in this paragraph, put
$\mathcal A_{N,B}=\{(N-p)/B:p<N,\ p\text{ prime},\ B\mid N-p\}$.
Let $\Phi_s$ be its positively marked average sifted at $R^{1/s}$.
Write $[l,r]^k_\uparrow$ for $k$ increasingly ordered exponents in
$[l,r]$.  For a positive integer $n$, let
\[
 Q_d(y;n)=\one_{d\mid n}\,
                  \one_{P^-(n/d)\ge y}.
\]
All additional primes in the following expressions are distinct;
collisions with the marked prefix, and square factors, are removed
with their fixed power-saving errors.  Every expression is then
averaged with the same nonnegative marked weights.
The first four terms are
\[
 \Gamma_1=4\Phi_{s_0}+\Phi_{\phi_1},\qquad
 \Gamma_2,\Gamma_3,\Gamma_4
   =\sum_{R^{a_0}\le q<R^c}Q_q(R^{a_0};n),
 \quad c=b,a_2,a_3,
\]
where the last expressions include the average over $n\in\mathcal A_{N,B}$.
For $\Gamma_5$ use the full region $D_5$ in
\eqref{sv:complete-fifth-region}; for $\Gamma_6,\Gamma_7,\Gamma_8$
use the other three pair regions above. The summand is
$Q_{q_1q_2}(R^{a_0};n)$ for $i=5,6$, and
$Q_{q_1q_2}(q_1;n)$ for $i=7,8$, with
$q_1=R^u<q_2=R^v$.
For $9\le i\le21$ the ordered exponent domains are
\begin{equation}\label{sv:switch-domains}
\begin{array}{c|l@{\qquad}c|l}
9 &[a_1,a_3]^3_\uparrow
 &10 &[a_1,a_2]^2_\uparrow[a_2,b]\\
11 &[a_1,a_2][a_2,a_3]^2_\uparrow
 &12 &[a_0,a_1]^2_\uparrow[a_3,b]\\
13 &[a_0,a_1][a_1,a_2][a_2,b]
 &14 &[a_0,a_1][a_2,b]^2_\uparrow\\
15 &[a_1,a_2][a_2,a_3][a_3,b]
 &16 &[a_2,a_3]^4_\uparrow\\
17 &[a_2,a_3]^3_\uparrow[a_3,b]
 &18 &[a_2,a_3]^2_\uparrow[a_3,b]^2_\uparrow\\
19 &[a_1,a_2][a_3,b]^3_\uparrow
 &20 &[a_2,a_3][a_3,b]^4_\uparrow\\
21 &[a_3,b]^6_\uparrow.&&
\end{array}
\end{equation}
For a domain with $r$ primes, the summand is
$Q_{q_1\cdots q_r}(q_{r-1};n)$.  Juxtaposition in
\eqref{sv:switch-domains} denotes independent successive bands,
with ordering only within each indicated ordered block.
The inequality to be used is
\begin{equation}\label{sv:twenty-one-weight}
 5\Phi_s\le\Gamma_1-\Gamma_2-\Gamma_3-\Gamma_4
                       +\sum_{i=5}^{21}\Gamma_i+o(\Theta_Q).
\end{equation}
It is the second weighted inequality of \cite[Lemma~4.2]{Wu2004};
the following verification fixes its complete ranges and multiplicities.

For a squarefree $R^{a_0}$-rough quotient let $A,B',C,D$ count
its prime factors in the four bands
$[a_0,a_1),[a_1,a_2),[a_2,a_3),[a_3,b)$, respectively.
(The symbol $B'$ in this verification is a band count, not the
marked prefix $B$.)  Pointwise the first four terms are
\[
 \Gamma_1=4+\one_{A=0},\qquad
 (\Gamma_2,\Gamma_3,\Gamma_4)
       =(A+B'+C+D,A+B',A+B'+C).
\]
The remaining terms have the following values, where
$x_+=\max(x,0)$:
\begin{equation}\label{sv:band-counts}
\begin{array}{c|l@{\quad}c|l}
5&\binom{A+B'}2&6&AC\\
7&(A-1)_+&8&\one_{A\ge1}B'\\
9&\one_{A=0}(B'+C-2)_+
 &10&\one_{A=0,B'\ge2}(C+D)\\
11&\one_{A=0,B'=1}(C-1)_+
 &12&\one_{A\ge2}D\\
13&\one_{A=1,B'\ge1}(C+D)
 &14&\one_{A=1,B'=0}(C+D-1)_+\\
15&\one_{A=0,B'=1,C\ge1}D
 &16&\one_{A=B'=0}(C-3)_+\\
17&\one_{A=B'=0,C\ge3}D
 &18&\one_{A=B'=0,C=2}(D-1)_+\\
19&\one_{A=0,B'=1,C=0}(D-2)_+
 &20&\one_{A=B'=0,C=1}(D-3)_+\\
21&\one_{A=B'=C=0}(D-5)_+.&&
\end{array}
\end{equation}
Indeed a quotient cutoff at the next-to-last chosen prime forces
all preceding chosen primes to be the smallest factors of $n$.
This gives each entry of \eqref{sv:band-counts}.  The signed
right-hand side in \eqref{sv:twenty-one-weight} is, respectively,
\[
\begin{array}{c|l}
A\ge2&(A+B'-2)(A+B'-3)/2+(A-2)C\\
A=1,\ B'\ge1&(B'-1)(B'-2)/2\\
A=1,\ B'=0&(1-C-D)_+\\
A=0,\ B'\ge2&(B'-2)(B'-3)/2\\
A=0,\ B'=1,\ C\ge1&0\\
A=0,\ B'=1,\ C=0&(2-D)_+\\
A=B'=0,\ C\ge3&0\\
A=B'=0,\ C=2,1,0&(1-D)_+,\ (3-D)_+,\ (5-D)_+.
\end{array}
\]
All are nonnegative, and at $A=B'=C=D=0$ the value is five.
The left side is five times the indicator of that last case.
Outside the $R^{a_0}$-rough support every summand vanishes.
This proves the pointwise inequality.  On a fixed positive exponent
box the omitted square and prefix-intersection terms have
$N^{1-\epsilon+o(1)}$ total mass for some fixed $\epsilon>0$,
which gives the stated error after the finite marked averaging.

\paragraph{Termwise estimates and the ordinary base.}
The two components of $\Gamma_1$ have upper multipliers
$4A_+(s_0)+A_+(\phi_1)$.  The three negative terms have lower
multipliers
\begin{equation}\label{sv:negative-singleton-base}
 \mathcal B(c)=\int_{a_0}^{c}
   \frac{A_-(s_0(1-u))}{u(1-u)}\,du,
 \qquad c=b,a_2,a_3.
\end{equation}
The support conditions put $s_0(1-u)$ in $(2,4]$ on these intervals,
so $A_-(s_0(1-u))=\log(s_0(1-u)-1)$.
The merged singleton and pair coefficients satisfy the complete
square and selected-complement tests of the parent state.  In
particular $D_5^+$, $D_6$, $D_7$, $D_8$ have the ordinary upper
kernel $(1-u-v)^{-1}$ against $du\,dv/(uv)$.
For a child argument below $2$, comparison with a smaller sifting
cutoff gives this same kernel, since $A_+=1$ up to $3$.
On
$D_5^-=\{a_0<u<v<a_2:v+2u\ge1\}$ the ordinary pair is not used.
The positive product estimate gives the valid kernel
$Kc_\omega s_0/4$, because its rough argument is at least
$s_0(2-2a_2)>2$.  Define
\begin{align}\label{sv:pair-base}
 \mathcal U_5&=
  \iint_{D_5^+}\frac{du\,dv}{uv(1-u-v)}
   +\frac{Kc_\omega s_0}{4}\iint_{D_5^-}\frac{du\,dv}{uv},\notag\\
 \mathcal U_i&=\iint_{D_i}\frac{du\,dv}{uv(1-u-v)}
                  \qquad(i=6,7,8).
\end{align}
Thus the unsupported corner is replaced by a positive upper
estimate; its ordinary mass or feedback is not retained as well.

For a switched domain $\mathcal D_i$ with $r$ primes, let
$w=u_{r-1}$ and $U=u_1+\cdots+u_r$.  Its normalized unsifted
mass at the actual parent ratio $\alpha$ is
\begin{equation}\label{sv:full-switched-mass}
 \mathcal M_i(\alpha)=
 \int_{\substack{\boldsymbol u\in\mathcal D_i\\U\le\alpha-w}}
 \frac{\omega((\alpha-U)/w)}{w\prod_j u_j}\,d\boldsymbol u
 +\int_{\substack{\boldsymbol u\in\mathcal D_i\\U=\alpha}}
                     \frac{du_1\cdots du_{r-1}}{\prod_j u_j}.
\end{equation}
In the surface integral one coordinate is eliminated with derivative
one.  The surface counts quotient $1$; it is separate from the
ordinary Buchstab term.  Prime summation and the rough-number
formula on fixed positive boxes give \eqref{sv:full-switched-mass}.
The upper count in parent units is $K\mathcal M_i(\alpha)/4$.

The following explicit bounds are uniform in $\alpha\ge2$.
Put $\ell_{ij}=\log(\phi_i/\phi_j)$,
$\ell_{01}=\log(s_0/\phi_1)$ and $\ell_{3b}=\log(\phi_3/s)$.
For $i=9,\ldots,15$ take $C_i=\max\{c_\omega,1/\max(1,m_i)\}$,
where
\[
\begin{array}{c|l@{\quad}c|l}
i&m_i&i&m_i\\ \hline
9&2\phi_3-3&10&2\phi_2-2-\phi_2/s\\
11&2\phi_3-\phi_3/\phi_2-2
 &12&2\phi_1-2-\phi_1/s\\
13&2\phi_2-\phi_2/\phi_1-1-\phi_2/s
 &14&2s-s/\phi_1-2\\
15&2\phi_3-\phi_3/\phi_2-1-\phi_3/s.&&
\end{array}
\]
These are lower bounds for the rough arguments on the respective
domains.  Equation~\eqref{sv:omega-envelope} gives
\begin{align}\label{sv:switch-mass-bounds}
 \widehat M_9&=C_9\{(\phi_1+\phi_3)\ell_{13}-2(\phi_1-\phi_3)\},\notag\\
 \widehat M_{10}&=C_{10}\{\phi_1-\phi_2(1+\ell_{12})\}
                                          \log(\phi_2/s),\notag\\
 \widehat M_{11}&=C_{11}\ell_{12}\{\phi_2(\ell_{23}-1)+\phi_3\},\notag\\
 \widehat M_{12}&=C_{12}\{s_0-\phi_1(1+\ell_{01})\}\ell_{3b},\notag\\
 \widehat M_{13}&=C_{13}\ell_{01}(\phi_1-\phi_2)\log(\phi_2/s),\notag\\
 \widehat M_{14}&=C_{14}\ell_{01}
                       \{\phi_2(\log(\phi_2/s)-1)+s\},\notag\\
 \widehat M_{15}&=C_{15}\ell_{12}(\phi_2-\phi_3)\ell_{3b},\notag\\
 \widehat M_{16}&=\phi_2\ell_{23}^4/24,
 &\widehat M_{17}&=\phi_2\ell_{23}^3\ell_{3b}/6,\notag\\
 \widehat M_{18}&=\phi_3\ell_{23}^2\ell_{3b}^2/4,
 &\widehat M_{19}&=\phi_3\ell_{12}\ell_{3b}^3/6,\notag\\
 \widehat M_{20}&=\phi_3\ell_{23}\ell_{3b}^4/24
                         +\phi_2\ell_{3b}^4/24,\notag\\
 \widehat M_{21}&=\phi_3\ell_{3b}^6/720+\phi_3\ell_{3b}^5/120.
 \end{align}
For $i\le19$ there is no unit surface because the sum of specified
exponents is at most $4/s<2\le\alpha$.
For $i=20$, eliminate the first variable on the surface and use
$1/u_1\le\phi_2$; the other four variables are ordered in
$[a_3,b]$.  This gives the second summand of $\widehat M_{20}$.
For $i=21$, eliminate the last variable, use $1/u_6\le\phi_3$,
and integrate the other five ordered variables.  The other
summands in \eqref{sv:switch-mass-bounds} are obtained by
integrating $1/(w\prod u_j)$ over their displayed bands, with
$\omega\le C_i$ for $9\le i\le15$ and $\omega\le1$
otherwise.  Hence $\mathcal M_i(\alpha)\le\widehat M_i$.

The ordinary base is now the explicit functional
\begin{equation}\label{sv:two-step-base}
 J_2^0=\frac15\left\{
 5A_+(s)-4A_+(s_0)-A_+(\phi_1)
 +\mathcal B(b)+\mathcal B(a_2)+\mathcal B(a_3)
 -\sum_{i=5}^8\mathcal U_i
 -\frac K4\sum_{i=9}^{21}\widehat M_i\right\}.
\end{equation}
Insert the upper estimates for positive terms and lower estimates
for negative terms in \eqref{sv:twenty-one-weight}.
The defects in the two components of $\Gamma_1$ subtract
$4H(s_0)+H(\phi_1)$ from its upper estimate.  The three singleton
lower defects subtract their integrals from the negative terms.
The four good-pair upper defects subtract the pair integrals from
their positive upper estimates.  After division by $5$ these are
exactly \eqref{sv:two-step-feedback}.  Therefore
\begin{equation}\label{sv:two-step-admissibility}
 \Phi_s\le\{A_+(s)-J_2^0-\mathcal P_2[H,h]+o(1)\}\Theta_Q,
\end{equation}
where $\mathcal P_2[H,h]$ is the expression in
\eqref{sv:two-step-feedback}.  This proves admissibility of a
finite two-step action from its mathematical base and positively
weighted preceding-generation defects.

Sharper upper bounds may replace the positive costs
in~\eqref{sv:pair-base} and~\eqref{sv:full-switched-mass}.
Such a bound must cover the same labelled family for every admissible
$\alpha\ge2$, including its unit surfaces. Replace that cost once
in~\eqref{sv:two-step-base}, and bound the resulting explicit
functional in the required direction to obtain the corresponding transition
constant. This preserves the sign argument for the two-step
inequality.

\subsection{A finite seed on the admitted marked families}

Put
\begin{equation}\label{sv:seed-parameters}
 S=\frac{11}{5},\qquad s_0=\frac{227}{50},\qquad
 (\phi_1,\phi_2,\phi_3)
 =\left(\frac{353}{100},\frac{29}{10},\frac{61}{25}\right),
 \qquad K_{\rm seed}=\frac{681}{100}.
\end{equation}
Here $K<K_{\rm seed}$, so Proposition~\ref{prop:current-positive-product}
also supplies the positive upper estimate with $K_{\rm seed}$.
Let $\mathcal J_{\rm seed}$ be the right-hand side of
\eqref{sv:two-step-base} at the parameters
\eqref{sv:seed-parameters}, with $s=S$ and $K$ replaced by
$K_{\rm seed}$. In particular its fifth cost is the full split cost
\eqref{sv:pair-base}, and the two unit-surface terms in
\eqref{sv:switch-mass-bounds} are included.

\begin{lemma}[Finite seed]\label{sv:finite-seed}
On an admitted marked family with $\alpha\ge2$, assume additionally
$L/S<\chi$ in the enhanced case. Then
\[
 \Phi_S\le\left(1-\frac{2327}{100000}+o(1)\right)\Theta_Q.
\]
The same upper defect is available at every $1\le w\le S$ whenever
the state admits the comparison cutoff $N^{L/S}$. The conclusion is
uniform over each fixed finite collection of such families, and follows
from four weighted steps after an initial seed of $1/50$.
\end{lemma}

\begin{proof}
All primes appended in the twenty-one-term inequality have exponent
at most $L/S$. They therefore lie below the retained ceiling, so the
largest specified slots, selected subsets and any third-slot bound
remain valid. The singleton residual square is $2u<1$. The pair squares
are $2v<1$ and $v+2u<1$; they hold on $D_5^+,D_6,D_7,D_8$.
Every new complete square is obtained by multiplying the corresponding
residual square by $B$. Once the selected old slots have occurred,
the new selected-complement square is the same full square. The
ordinary residual atoms are evaluated at their actual level
$Q/(Bq)$ or $Q/(Bq_1q_2)$. The remaining part $D_5^-$ is retained
with its positive cost, using no signed distribution claim there.
In the classical case the complete ordinary modulus is below $Q$,
and the absolute below-half remainder gives the same conservative
functional. The switched terms require only $\alpha\ge2$ and fixed
positive prime cutoffs.

The directed enclosure of the explicit functional is
\begin{equation}\label{sv:seed-base-certificate}
 \mathcal J_{\rm seed}\ge
 \frac{11672890052052913}{576460752303423488}
 >\frac{20249}{10^6}>\frac1{50}.
\end{equation}
The finite certificate for this inequality is specified in the
ancillary data under \texttt{seed\_certificate}. Its ordinary
integrals are those in \eqref{sv:negative-singleton-base} and
\eqref{sv:pair-base}; its switched masses are the explicit upper
bounds in \eqref{sv:switch-mass-bounds}. The ordinary double integral
is first integrated in one coordinate. The resulting one-dimensional
logarithmic integrals are enclosed by rational Darboux intervals and
the logarithm remainder of \eqref{num:log-enclosure}. Upper endpoints
are used for positive costs and lower endpoints for negative costs in
the definition of $\mathcal J_{\rm seed}$. Thus
\eqref{sv:seed-base-certificate} proves the initial defect $J_0=1/50$
before any later array recurrence is used.

Suppose an upper defect $J>0$ on $[1,S]$ has already been proved on
these admitted families. Positivity and cutoff comparison give
\[
 H_J(w)=J\one_{[1,S]}(w).
\]
The lower Buchstab identity at the same total level gives the
previously proved option
\begin{equation}\label{sv:seed-lower-child}
 h_J(w)=J\log\frac{S}{w-1}\,
                    \one_{[2,1+S]}(w).
\end{equation}
The normalization is independent of $w$, and $A_+(w)=1$ for
$1\le w\le S$. To check the comparison inside
\eqref{sv:seed-lower-child}, a Buchstab child indexed by
$v\in[w-1,S]$ has appended exponent $L/(v+1)$ and residual
exponent $Lv/(v+1)$. Its seed comparison exponent satisfies
\[
 \frac{Lv}{S(v+1)}\le\frac{L}{v+1}.
\]
It is therefore below that appended prime and the preceding stable
ceiling. Its new square holds because $v\ge1$. The lower child
need not itself admit argument $S$; the seed is applied to its
upper Buchstab children.

Insert $H_J,h_J$ into \eqref{sv:two-step-feedback}. Since
$s_0,\phi_1>S$, the unchanged-prefix defects vanish. Write
$u_*=67/227$, $b=1/S$, and set
\[
 E_i=\begin{cases}
 D_i\cap\{u+v\ge117/227\},&i=5^+,6,\\
 D_i\cap\{v\ge1-16u/5\},&i=7,8.
 \end{cases}
\]
The feedback is exactly $\lambda_{\rm seed} J$, where
\begin{align}\label{sv:seed-feedback-integral}
 \lambda_{\rm seed}=\frac15\bigg\{&
 \sum_{c\in\{b,a_2,a_3\}}
 \int_{u_*}^{c}
 \frac{\log\{S/[s_0(1-u)-1]\}}{u(1-u)}\,du\notag\\
 &+\sum_{i=5^+,6,7,8}\iint_{E_i}
                     \frac{du\,dv}{uv(1-u-v)}\bigg\}.
\end{align}
For $i=5^+,6$ the support condition is
$s_0(1-u-v)\le S$; for $i=7,8$ it is
$(1-u-v)/u\le S$. All retained pair-child arguments exceed $1$,
and the singleton lower arguments exceed $2$. There is no
feedback from $D_5^-$.

Every boundary in \eqref{sv:seed-feedback-integral} is rational
and affine. For a vertical slice $l(u)<v<r(u)$ the exact pair
integral is
\[
 \frac1{u(1-u)}
 \log\frac{r(u)(1-u-l(u))}{l(u)(1-u-r(u))}.
\]
For $E_{5^+}$ the successive $u$ intervals have endpoints
$50/227,117/454,19/58,1/3$. The domain $E_6$ is its entire
rectangle. The nonempty support in $E_7$ starts at $1265/5648$
and changes its lower boundary at $5/21$; the lower boundary of
$E_8$ changes at $1265/5648$. Rational Darboux enclosures of
these pieces give
\begin{equation}\label{sv:seed-feedback-certificate}
 \lambda_{\rm seed}\ge
 \frac{75057365825327405}{576460752303423488}
 >\frac{13}{100}.
\end{equation}
The exact endpoints, subdivisions and source correspondence are
specified by the same \texttt{seed\_certificate} data as
\eqref{sv:seed-base-certificate}.

Equations \eqref{sv:two-step-admissibility},
\eqref{sv:seed-base-certificate} and
\eqref{sv:seed-feedback-certificate} prove
\[
 J_{n+1}=\frac{20249}{10^6}+\frac{13}{100}J_n.
\]
At each step \eqref{sv:seed-lower-child} is derived from $J_n$
before the weighted inequality is applied. Starting at $J_0=1/50$
the four exact values are
\[
 \frac{22849}{10^6},\quad
 \frac{2321937}{10^8},\quad
 \frac{232675181}{10^{10}},\quad
 \frac{23273777353}{10^{12}}>\frac{2327}{100000}.
\]
There are finitely many appended-prime operations. Their exponents
remain positive after all compact choices have been fixed, so the
positive-product estimate and the distribution remainders apply at
each of them. The finitely many errors are $o(\Theta_Q)$; the strict
final surplus absorbs them in the stated order of limits.
\end{proof}

\paragraph{Use at the roots and descendants.}
For $U_3,U_4$, the specified prime has exponent $t>{a_{\rm s}}$, the residual
exponent is $L=v(t)-t$, and the upper argument is
$\max(2,L/{a_{\rm s}})$. On the enhanced range $t<57/224$, the largest-prime branch has
$2t<v(t)\le(1+t)/2$, so its stable ceiling can be taken below $t$
and arbitrarily close to it. Its physical cutoff exponent
$\min({a_{\rm s}},L/2)\le {a_{\rm s}}<t$ lies below that ceiling. Whenever the seed is used, its comparison exponent
$L/S$ is no larger than this physical cutoff. The classical tail
uses the below-half interface. For $U_{10},U_{11}$ the same argument
uses the least specified exponent $w$ and physical cutoff
$\min(w,L/2)$; the original selected subsets and first three slots
are already fixed. For the lower pair terms $L_5,L_6$, the seed
occurs only inside \eqref{sv:seed-lower-child}. The first appended
prime is at most $N^{a_{\rm s}}$, and every seed comparison prime is below
it by the preceding displayed inequality. These comparisons also
preserve a third-slot bound equal to ${a_{\rm s}}$.

More generally, every finite action in
\eqref{sv:one-step-operator} or \eqref{sv:two-step-feedback} retains
the total $Q$, every old square and the stable ceiling. New
singletons have $2u<1$; new pairs in the ordinary part have
$2v<1$ and $v+2u<1$. Their physical cutoffs are no larger than
the least newly appended prime. Equation~\eqref{sv:child-state}
shows that their ratios remain at least $2$. Induction therefore
keeps the seed and its Buchstab consequences within the admitted
families. No seed defect is inserted directly into an empty-prefix
lower start.

The finite induction uses the preceding generation on every
right-hand side. At each state take the maximum of alternative
estimates; when a child argument crosses a knot, use the minimum
over all adjoining closed cells. Each step therefore gives a bound
for the same positive marked counts. Only finitely many steps
are used.

Here is an explicit common approximation bound.  The positive operator
masses satisfy
\begin{align*}
 \|\mathcal B\|&\le\log5<2,\\
 \|\mathcal W_1\|&\le1+\tfrac12\log4<2,\\
 \|\mathcal W_2\|&\le1+\tfrac35\log4
                 +\tfrac45\,125(3/10)^2<12.
\end{align*}
For the last estimate $1/5\le u,v\le1/2$ and
$1-u-v\ge u\ge1/5$, so the pair kernel is at most $125$ on
each of four regions.  The unchanged-prefix coefficients sum to one.
The switched positive costs belong to the base of the action and
are not additional feedback edges.
If each local oriented approximation error is at most $\delta$,
then
\begin{equation}\label{sv:feedback-error}
 e_{j+1}\le\delta+12e_j,\qquad
 e_j\le\delta\frac{12^{j+1}-1}{11}.
\end{equation}
The absolute coefficients of the eleven-term functional sum to $14$
before its factor $1/4$, giving root norm $7/2$.

At depth at most $440$, every residual level is at least
$\lambda_*101^{-440}>0$, and an appended prime exponent is at
least this number divided by $6$.  Recursively refine appended
exponents on each fixed positive parent, using an error within
$\delta$ times its normalized mass.  Proper cap and knot walls are
null; structural clamps and inherit identities are retained.
The denominators are bounded away from zero.  Equation
\eqref{sv:feedback-error} then gives one finite set of queries with
arbitrarily small total error.  Only after fixing that set does $N$
increase.  A positive-mass vertex has a positive limiting
$\Theta_Q/H_N$, so the absolute distribution errors become
$o(\Theta_Q)$; zero-mass pieces are handled by absolute errors.

\section{Endpoint estimates and the order of limits}\label{app:cross-review}

The sharp-prime lemma (Lemma~\ref{lem:current-sharp-prime}) applies to
the same coefficient constructed above.  We record the absolute bounds
that allow it to be used for signed lower weights.  If
$|\rho_N^{\rm low}(q)|\le C$, its squarefree support is $q\le N^\theta$,
$\theta<1$, and $\delta=(\log N)^{-B}$, then the endpoint contribution
is bounded by
\begin{equation}\label{sv:sharp-boundary}
 \sum_{\substack{p\text{ prime}\\0<|p-N|\le\delta N}}
          \sum_{q\mid N-p}|\rho_N^{\rm low}(q)|
 \le2C\sum_{1\le n\le\lceil\delta N\rceil}\tau(n)
 \ll\delta N\log N+\log N.
\end{equation}
The difference $p-N=0$ is impossible for even $N>2$.
The exact-coprime mean boundary has the same bound because
$\sum_{q\le N^\theta}1/\varphi(q)\ll\log N$.
The smooth lower trim below $2N^{1-\kappa_0}$ costs
$N^{1-\kappa_0+o(1)}$ in the progression sum and
$O(N^{1-\kappa_0}\log N)$ in the mean.  These have fixed power
savings.  Choosing $B>A+3$ makes all boundary terms smaller than the
prescribed logarithmic error.

For the exact-coprime prime mean, let $\omega_{\mathrm{pf}}(q)$ be
the number of distinct prime divisors of $q$.  Every such divisor is
below $N$, and squarefreeness gives
\begin{equation}\label{sv:coprime-prime-error}
 \sum_q\frac{|\rho_N^{\rm low}(q)|\omega_{\mathrm{pf}}(q)}{\varphi(q)}
 \ll\sum_{\ell\le N^\theta}\frac1{\ell-1}
                  \sum_{e\le N^\theta/\ell}\frac1{\varphi(e)}
 \ll\log^2N.
\end{equation}
The sum over $\ell$ is over primes; summing over all integers already
gives the stated crude bound.  Removing the $O(\log N)$ prime divisors
of $N$ from the original sequence contributes $N^{o(1)}$ by the
divisor bound.  These estimates justify using the sharp original
prime count in \eqref{sv:physical-lower}, without obtaining a prime
lower bound from the positive-product upper constant $K$.

The upper and lower count estimates follow the parameter order in
Section~\ref{sec:lower-limits}. The fixed moment orders precede the
choice of $\kappa$ in~\eqref{coef:parameter-order}, and the boundary
exponent satisfies $B>A+3$, as above.
